\documentclass[../../../include/open-logic-section]{subfiles}

\begin{document}

\olfileid[tr]{his}{set}{cantorplane}
\olsection{Cantor'da Doğru ve Düzlem}

Schr\"oder--Bernstein ispatını çevreleyen koşulların bazıları,
\olref[his][set][pathology]{sec}'te ele aldığımız tarihle bağlantılıdır.
Cantor'un 1877'de bir karenin üzerindeki noktalarla bir kenarının üzerindeki
noktaların tam olarak aynı sayıda olduğunu ispatladığını anımsayın. Burada
onun (ilk giriştiği) ispatı sunacağız.

$\unitline$, birim doğru parçası, yani $[0,1]$ noktaları kümesi olsun.
$\unitsquare$, birim kare, yani $\unitline \times \unitline$ noktaları kümesi
olsun. Bu terimlerle Cantor $\cardeq{\unitline}{\unitsquare}$ olduğunu
ispatladı. Dedekind'e, özünde aşağıdaki argümanı içeren bir not yazdı.

\begin{thm}\ollabel{thm:cantorplane}
$\cardeq{\unitline}{\unitsquare}$
\end{thm}

\begin{proof}[İspat: birinci kısım.] 
$a, b \in \unitline$ sabitleyin. Bunları ikili gösterimle yazın; böylece
$0$'lar ile $1$'lerden oluşan $a_1$, $a_2$, \dots{} ve $b_1$, $b_2$,
\dots{} sonsuz dizilerini şöyle elde ederiz:
\begin{align*}
a &= 0.a_1a_2a_3a_4\dots\\
b &= 0.b_1b_2b_3b_4\dots
\intertext{Şimdi şu biçimde verilen $f \colon \unitsquare \to \unitline$
fonksiyonunu ele alın:} 
f(a, b) & = 0.a_1b_1a_2b_2a_3b_3a_4b_4\dots
\end{align*}
$f$ bir !!a{injection}{}dur; çünkü $f(a, b) = f(c,d)$ ise $a_n = c_n$ ve
$b_n = d_n$ olur; bu, her $n \in \Nat$ için geçerlidir, dolayısıyla $a = c$
ve $b = d$'dir.
\end{proof}

Ne yazık ki Dedekind'in Cantor'a belirttiği gibi bu, özgün soruyu yanıtlamaz.
$0.\dot{1}\dot{0} = 0.1010101010\ldots$ sayısını ele alın.
$f(a,b) = 0.\dot{1}\dot{0}$ olmasına ihtiyacımız vardır; burada:
\begin{align*}
a&= 0.\dot{1}\dot{1} = 0.111111\ldots\\
b&= 0
\end{align*}
Fakat $a = 0.\dot{1}\dot{1} = 1$'dir. Öyleyse ``$a$ ile $b$'yi ikili
gösterimle yazın'' dediğimizde \emph{hangi} gösterimi kullanacağımızı
seçmeliyiz; $f$ bir \emph{fonksiyon} olacağı için olası iki gösterimden
yalnızca \emph{birini} kullanabiliriz. Ama örneğin sonlu gösterimi kullanıp
$a$'yı ``$1.000\ldots$'' diye yazarsak hiçbir $\tuple{a, b}$ ikilisi için
$f(a, b) = 0.\dot{1}\dot{0}$ olmaz.

Özetlersek Dedekind, belirli yinelenen ikili açılımların mümkün olması
nedeniyle Cantor'un $f$ fonksiyonunun bir !!a{injection} olduğunu, fakat bir
!!a{surjection} \emph{olmadığını} belirtti. Dolayısıyla Cantor yalnızca
$\cardle{\unitsquare}{\unitline}$ olduğunu göstermiştir;
$\cardeq{\unitsquare}{\unitline}$ olduğunu \emph{değil}.

Cantor hemen hemen anında Dedekind'e yanıt yazarak ispatın özünde şöyle
tamamlanabileceğini ileri sürdü:

\begin{proof}[İspat: tamamlanmış biçimi.] 
Demek ki $\cardle{\unitsquare}{\unitline}$ olduğunu gösterdik. Ancak
$\unitline$'dan $\unitsquare$'e açıkça bir !!a{injection} vardır: doğru
parçasını karenin bir kenarı boyunca yatırmak yeterlidir. Öyleyse
$\cardle{\unitline}{\unitsquare}$ ve
$\cardle{\unitsquare}{\unitline}$ olur. Schr\"{o}der--Bernstein gereği
(\olref[sfr][siz][sb]{thm:schroder-bernstein}),
$\cardeq{\unitline}{\unitsquare}$.
\end{proof}

Ama Cantor elbette son satırı bu terimlerle tamamlayamazdı; çünkü
Schr\"{o}der--Bernstein Teoremi henüz ispatlanmamıştı. Gerçekten de Cantor
daha sonra bunu genel bir varsayım olarak formüle edecek olsa da doyurucu bir
ispat ancak 1897'de verildi. (Dolayısıyla Cantor 1877'nin ilerleyen
aylarında, Schr\"{o}der--Bernstein üzerinden gitmeyen farklı bir
\olref{thm:cantorplane} ispatı sundu.)

\end{document}
